\documentclass[12pt,a4paper]{article}
\usepackage[utf8]{inputenc}
\usepackage[T1]{fontenc}
\usepackage[brazil]{babel}
\usepackage{geometry}
\geometry{top=2.5cm,bottom=2.5cm,left=2.5cm,right=2.5cm}
\usepackage{amsmath,amsfonts,amssymb}
\usepackage{graphicx}
\usepackage{booktabs}
\usepackage{xcolor}
\usepackage{hyperref}
\usepackage{listings}
\usepackage{tikz}
\usetikzlibrary{shapes.geometric,arrows.meta,positioning,calc,fit,backgrounds}

\hypersetup{
    colorlinks=true,
    linkcolor=blue!70!black,
    urlcolor=blue!70!black,
    citecolor=blue!70!black
}

\lstset{
    basicstyle=\ttfamily\small,
    keywordstyle=\color{blue!80!black}\bfseries,
    stringstyle=\color{green!50!black},
    commentstyle=\color{gray}\itshape,
    frame=single,
    breaklines=true,
    numbers=left,
    numberstyle=\tiny\color{gray},
    backgroundcolor=\color{black!3},
    showspaces=false,
    showstringspaces=false
}

\title{\textbf{Descrição dos Padrões de Software Utilizados}\\
\large Sistema de Monitoramento Hospitalar Baseado em IoT e Mensageria (Arquitetura sem Exchanges)}
\author{\textbf{Equipe de Arquitetura de Software}}
\date{\today}

\begin{document}

\maketitle

\begin{abstract}
Este documento apresenta a especificação e justificativa técnica dos padrões de projeto (\textit{Design Patterns}) empregados no desenvolvimento do Sistema de Monitoramento Hospitalar IoT. São discutidas as motivações para a substituição do padrão \textit{Facade} pelo padrão \textbf{Mediator}, a eliminação completa das \textit{Exchanges} complexas no RabbitMQ (transferindo a inteligência de roteamento clínico para o próprio Mediator), bem como a aplicação dos padrões \textit{Template Method} e \textit{Simple Factory}. Inclui-se o Diagrama de Classes UML que fundamenta a estrutura orientada a objetos da solução.
\end{abstract}

\tableofcontents
\newpage

\section{Introdução e Contexto da Refatoração}
O sistema desenvolvido tem como objetivo coletar, classificar e distribuir telemetrias de sinais vitais de pacientes hospitalares (frequência cardíaca, pressão arterial, saturação de oxigênio $O_2$ e temperatura corporal) para diferentes setores clínicos: Posto de Enfermagem, Painel do Médico Especialista e Aplicativo Móvel de Emergência.

\subsection{A Decisão de Eliminar as Exchanges do RabbitMQ}
Nas versões preliminares do projeto, o sistema dependia de um \textit{Topic Exchange} no RabbitMQ com chaves hierárquicas complexas (\texttt{<paciente>.<sinal>.<prioridade>}) para realizar o roteamento nos clientes. Em discussão arquitetural com o professor, identificou-se que essa abordagem transferia indevidamente regras do domínio da aplicação para a infraestrutura de mensageria, além de dificultar o desacoplamento e a rastreabilidade dos nós sensores.

Optou-se por \textbf{eliminar completamente qualquer declaração de Exchange} (Topic, Direct ou Fanout). No novo modelo:
\begin{enumerate}
    \item Cada nó sensor publica diretamente em sua própria fila ponto-a-ponto no RabbitMQ (\texttt{fila.cardiaco}, \texttt{fila.pressao}, \texttt{fila.saturacao}, \texttt{fila.temperatura}).
    \item O componente \textbf{Mediator} atua como o cérebro do sistema: consome essas quatro filas, executa a análise de dados e aplica as regras de triagem hospitalar.
    \item O próprio \textbf{Mediator} entrega os dados nas filas de saída de cada monitor (\texttt{fila.enfermagem}, \texttt{fila.medico.<id>}, \texttt{fila.alerta}).
\end{enumerate}

---

\section{Padrões de Software Empregados}

\subsection{Padrão Comportamental: Mediator}

\subsubsection{Motivação e Responsabilidades}
O padrão \textbf{Mediator} (GoF) define um objeto que encapsula e orquestra a forma como um conjunto de entidades interage, promovendo o acoplamento fraco ao evitar que as partes se comuniquem diretamente.

No sistema hospitalar, o \texttt{MonitorMediator} centraliza duas grandes responsabilidades que antes estavam desarticuladas:
\begin{enumerate}
    \item \textbf{Consumo Ponto-a-Ponto e Análise Clínica:} O mediador escuta de forma assíncrona as quatro filas individuais de sensores. Ao receber uma medição, utiliza verificações de tipo (\texttt{instanceof}) para acionar a rotina correta de interpretação fisiológica (\textit{Normal}, \textit{Atenção}, \textit{Crítico}).
    \item \textbf{Roteamento Inteligente sem Exchanges:} O mediador implementa a lógica de distribuição:
    \begin{itemize}
        \item \textbf{Enfermagem:} Recebe todas as mensagens de todos os leitos $\rightarrow$ o Mediator publica na \texttt{fila.enfermagem}.
        \item \textbf{Médico:} Recebe somente as leituras do paciente sob seus cuidados $\rightarrow$ o Mediator publica na fila dedicada \texttt{fila.medico.<idPaciente>}.
        \item \textbf{Móvel / Emergência:} Recebe apenas eventos com risco de vida $\rightarrow$ o Mediator publica na \texttt{fila.alerta} exclusivamente quando a prioridade clínica for \texttt{"critico"}.
    \end{itemize}
\end{enumerate}

\subsubsection{Por que substituir o Facade pelo Mediator?}
A Tabela~\ref{tab:facade_vs_mediator} sintetiza o comparativo arquitetural:

\begin{table}[h!]
\centering
\small
\begin{tabular}{p{4.2cm}p{5.5cm}p{5.5cm}}
\toprule
\textbf{Critério} & \textbf{Padrão Facade (Anterior)} & \textbf{Padrão Mediator (Atual)} \\
\midrule
\textbf{Papel Arquitetural} & Apenas mascarava métodos dos sensores dentro de um único processo em memória. & Intermediário autônomo entre processos produtores e processos consumidores. \\
\textbf{Infraestrutura} & Obrigava a usar Topic Exchange com casamento de wildcards na mensageria. & Permite usar filas simples diretas, eliminando totalmente exchanges. \\
\textbf{Autonomia dos Sensores} & Nula: o Facade instanciava e media todos os sensores em um loop monolítico. & Alta: cada sensor é um publicador independente conectado à sua fila. \\
\textbf{Localização da Regra} & Regras de negócio misturadas com loop de transmissão do broker. & Lógica de negócio e roteamento totalmente concentradas no Mediator. \\
\bottomrule
\end{tabular}
\caption{Comparação entre Facade e Mediator no sistema hospitalar.}
\label{tab:facade_vs_mediator}
\end{table}

\subsubsection{Discussão sobre o Padrão Visitor}
Durante os estudos preliminares, cogitou-se utilizar o padrão \textbf{Visitor} para substituir os testes de \texttt{instanceof} no mediador por \textit{double-dispatch}. No entanto, considerando que o domínio do sistema abrange estritamente quatro sinais vitais hospitalares padronizados, a equipe optou pela simplicidade do \textbf{Mediator} puro: o uso de \texttt{instanceof} é direto, não prejudica o desempenho e evita a criação de hierarquias de visitantes desnecessárias para um escopo acadêmico claro e conciso.

\subsection{Padrão Arquitetural: Mensageria Ponto-a-Ponto e Filas Dedicadas}

Com a eliminação da \textit{Topic Exchange}, o sistema opera exclusivamente com o mecanismo de \textbf{Filas Ponto-a-Ponto} no RabbitMQ:
\begin{itemize}
    \item \textbf{Filas de Entrada (Sensores $\rightarrow$ Mediator):}
    \begin{itemize}
        \item \texttt{fila.cardiaco}: buffer exclusivo para leituras de BPM.
        \item \texttt{fila.pressao}: buffer exclusivo para leituras sistólica/diastólica.
        \item \texttt{fila.saturacao}: buffer exclusivo para $SpO_2$.
        \item \texttt{fila.temperatura}: buffer exclusivo para temperatura em $^\circ\text{C}$.
    \end{itemize}
    \item \textbf{Filas de Saída (Mediator $\rightarrow$ Monitores):}
    \begin{itemize}
        \item \texttt{fila.enfermagem}: consumida pelo painel geral da enfermagem.
        \item \texttt{fila.medico.<id>}: consumida pelo terminal do médico do leito.
        \item \texttt{fila.alerta}: consumida pelo aplicativo móvel de emergência.
    \end{itemize}
\end{itemize}
Dessa forma, cada monitor consome de sua respectiva fila sem concorrência indevida de mensagens com outros setores (evitando que vire uma \textit{Work Queue} destrutiva acidental).

\subsection{Padrão Comportamental: Template Method}

Implementado na classe abstrata \texttt{PainelHospitalar}:
\begin{lstlisting}[language=Java]
public abstract class PainelHospitalar {
    protected abstract String getNomeFila();
    protected abstract void processarMensagem(String mensagem);
    protected abstract String getNomePainel();

    // Template Method invariante
    public final void iniciar() {
        // 1. Conecta ao RabbitMQ
        // 2. Declara a fila direta getNomeFila() (sem exchange)
        // 3. Registra DeliverCallback com processarMensagem()
        // 4. Inicia consumo assincrono
    }
}
\end{lstlisting}
O método \texttt{iniciar()} é declarado como \texttt{final}, garantindo que todas as subclasses (\texttt{PainelEnfermagem}, \texttt{PainelMedico}, \texttt{AppEmergencia}) compartilhem rigorosamente o mesmo fluxo de conexão e escuta assíncrona.

\subsection{Padrão Criacional: Simple Factory}

Implementado na classe \texttt{PainelFactory}, isolando a instanciação dos tipos de clientes hospitalares:
\begin{lstlisting}[language=Java]
public class PainelFactory {
    public static PainelHospitalar criarPainel(String tipo, String paciente) {
        switch (tipo.toLowerCase()) {
            case "enfermagem": return new PainelEnfermagem();
            case "medico":     return new PainelMedico(paciente);
            case "emergencia": return new AppEmergencia(paciente);
            default:           return null;
        }
    }
}
\end{lstlisting}

\newpage
\section{Diagrama de Classes da Solução}

A Figura~\ref{fig:diagrama_classes} ilustra a estrutura de classes orientada a objetos do projeto.

\begin{figure}[h!]
\centering
\begin{tikzpicture}[
    classbox/.style={draw=black!80, fill=blue!5, rounded corners=2pt, font=\footnotesize, text width=4.3cm},
    abstractbox/.style={draw=black!80, fill=green!5, rounded corners=2pt, font=\footnotesize, text width=4.3cm},
    arrow/.style={-{Latex[length=2.5mm]}, line width=0.8pt},
    inherit/.style={-{Triangle[open,length=3mm]}, line width=0.8pt},
    dashedinherit/.style={-{Triangle[open,length=3mm]}, line width=0.8pt, dashed}
]

% Interface Sensor
\node[classbox] (sensor) at (0, 7.5) {
    \begin{center}\textbf{<<interface>>}\\\textbf{Sensor}\end{center}
    \rule{4.3cm}{0.4pt}\\
    + medir(): double\\
    + getTipoSinal(): String\\
    + getDescricao(): String
};

% Classes Concretas de Sensores
\node[classbox, text width=3.2cm] (cardiaco) at (-4.5, 4.5) {
    \begin{center}\textbf{Cardiaco}\end{center}
    \rule{3.2cm}{0.4pt}\\
    - bpm: double\\
    \rule{3.2cm}{0.4pt}\\
    + medir(): double\\
    + getBpm(): double
};

\node[classbox, text width=3.2cm] (pressao) at (-1.2, 4.5) {
    \begin{center}\textbf{Pressao}\end{center}
    \rule{3.2cm}{0.4pt}\\
    - sistolica: double\\
    - diastolica: double\\
    \rule{3.2cm}{0.4pt}\\
    + medir(): double
};

\node[classbox, text width=3.2cm] (saturacao) at (2.1, 4.5) {
    \begin{center}\textbf{Saturacao}\end{center}
    \rule{3.2cm}{0.4pt}\\
    - spo2: double\\
    \rule{3.2cm}{0.4pt}\\
    + medir(): double
};

\node[classbox, text width=3.2cm] (temp) at (5.4, 4.5) {
    \begin{center}\textbf{Temperatura}\end{center}
    \rule{3.2cm}{0.4pt}\\
    - temp: double\\
    \rule{3.2cm}{0.4pt}\\
    + medir(): double
};

% Conexoes Sensores
\draw[dashedinherit] (cardiaco.north) -- ++(0,0.5) -| (sensor.south);
\draw[dashedinherit] (pressao.north) -- ++(0,0.5) -| (sensor.south);
\draw[dashedinherit] (saturacao.north) -- ++(0,0.5) -| (sensor.south);
\draw[dashedinherit] (temp.north) -- ++(0,0.5) -| (sensor.south);

% Interface SensorMediator
\node[classbox] (mediatorInt) at (7, 8) {
    \begin{center}\textbf{<<interface>>}\\\textbf{SensorMediator}\end{center}
    \rule{4.3cm}{0.4pt}\\
    + iniciar(): void\\
    + processar(pac, s): Resultado\\
    + analisar(s): Resultado
};

% Classe Concreta MonitorMediator
\node[classbox, text width=5cm] (mediatorImpl) at (7, 4.5) {
    \begin{center}\textbf{MonitorMediator}\end{center}
    \rule{5cm}{0.4pt}\\
    - channel: Channel\\
    \rule{5cm}{0.4pt}\\
    + iniciar(): void\\
    + processarMensagemFila(f, p): void\\
    + processar(pac, s): Resultado\\
    + analisar(s): Resultado
};
\draw[dashedinherit] (mediatorImpl.north) -- (mediatorInt.south);

% Hierarquia de Assinantes (Template Method)
\node[abstractbox, text width=5cm] (painelHosp) at (0, 1.2) {
    \begin{center}\textbf{<<abstract>>}\\\textbf{PainelHospitalar}\end{center}
    \rule{5cm}{0.4pt}\\
    \# getNomeFila()*: String\\
    \# processarMensagem(msg)*: void\\
    \# getNomePainel()*: String\\
    \rule{5cm}{0.4pt}\\
    + iniciar(): void
};

\node[classbox, text width=3.3cm] (enfermagem) at (-4, -1.8) {
    \begin{center}\textbf{PainelEnfermagem}\end{center}
    \rule{3.3cm}{0.4pt}\\
    \# getNomeFila(): "fila.enfermagem"\\
    \# processarMensagem()
};

\node[classbox, text width=3.3cm] (medico) at (0, -1.8) {
    \begin{center}\textbf{PainelMedico}\end{center}
    \rule{3.3cm}{0.4pt}\\
    - idPaciente: String\\
    \rule{3.3cm}{0.4pt}\\
    \# getNomeFila(): "fila.medico.<id>"\\
    \# processarMensagem()
};

\node[classbox, text width=3.3cm] (emergencia) at (4, -1.8) {
    \begin{center}\textbf{AppEmergencia}\end{center}
    \rule{3.3cm}{0.4pt}\\
    - idPaciente: String\\
    \rule{3.3cm}{0.4pt}\\
    \# getNomeFila(): "fila.alerta"\\
    \# processarMensagem()
};

\draw[inherit] (enfermagem.north) -- ++(0,0.5) -| (painelHosp.south);
\draw[inherit] (medico.north) -- (painelHosp.south);
\draw[inherit] (emergencia.north) -- ++(0,0.5) -| (painelHosp.south);

\end{tikzpicture}
\caption{Diagrama de Classes UML com Mediator, Filas Diretas e Template Method.}
\label{fig:diagrama_classes}
\end{figure}

\section{Conclusão}
A eliminação das \textit{Exchanges} complexas transferiu a lógica de roteamento para onde ela realmente pertence: no \textbf{MonitorMediator}. Dessa forma, a arquitetura tornou-se simultaneamente mais simples do ponto de vista de infraestrutura e muito mais expressiva no nível do software, atendendo com fidelidade o esboço de implantação validado com o professor.

\end{document}
